\begin{resumo}

Os organismos artificiais celulares são circuitos que encapsulam e controlam arquiteturas integradoras, capacitadas para a execução de todos os cálculos requeridos sobre o domínio discreto escolhido. Os autômatos celulares constituem uma alternativa para modelagem de sistemas físicos em domínio discreto. Neste trabalho, o fenômeno físico considerado como estudo de caso é a propagação de ondas sonoras através de um meio ideal e estacionário. O trabalho foi realizado em duas etapas, sendo a primeira voltada para ao desenvolvimento de um sistema celular autômato (CA) aplicável à emulação de ondas planas unidimensionais, para aplicação em circuitos reconfiguráveis, e a segunda voltada para a apresentação de propostas de implementação em chip. O sistema CA desenvolvido, denominado Camphs1D, é composto de duas células autômatas, baseadas em IP-cores de operadores aritméticos em 27 bits, capazes de operar todos os fenômenos básicos de onda em meios unidimensionais. As mencionadas células foram desenvolvidas a partir de duas estratégias distintas, sendo a primeira descrita diretamente em VHDL, no software Vivado, e a segunda descrita em C/C++, no software Vivado HLS, onde foram usados IP-cores de ponto flutuante de 32 bits da Xilinx. Foram utilizadas, para efeito de validação de resultados, simulações em \textit{Testbench} VHDL, simulações no software SDK (Xilinx), e implementações físicas integradas com um processador ARM, para acesso à memória DDR disponível em placa. Todas as validações atingiram seus objetivos, relacionados à comparações com os mesmos resultados obtidos por modelos de referência, desenvolvidos a partir da solução de d'Alembert para a equação unidimensional da onda.

 \vspace{\onelineskip}
    
 \noindent
 \textbf{Palavras-chaves}: autômatos celulares; malha de guias digitais de ondas; sistemas em chip.
\end{resumo}
